\chapter{Conclusiones y trabajos futuros}
\section{Conclusiones}

El presente trabajo tuvo como objetivo evaluar el desempeño del protocolo ANTop en redes satelitales de tipo DTN y compararlo con dos variantes de CGR: \textit{one-best route} y \textit{per-neighbor route}. Para ello se realizaron simulaciones sobre constelaciones satelitales de distinta escala y bajo diferentes niveles de fallas, analizando métricas clave como \textit{Arrival Time}, \textit{Elapsed Time}, \textit{Number of Hops} y \textit{Delivery Ratio}.

Los resultados obtenidos muestran que ANTop presenta el mejor desempeño global entre los algoritmos evaluados. En términos de latencia, logra mantener consistentemente los menores tiempos de llegada de paquetes, incluso a medida que aumenta el tamaño de la constelación. Asimismo, el número de saltos requerido para alcanzar el destino se mantiene bajo y con una variabilidad reducida, lo que indica que el algoritmo es capaz de encontrar rutas eficientes de forma consistente.

Desde el punto de vista computacional, ANTop demuestra un comportamiento estable y predecible. A diferencia de las variantes de CGR, cuyos tiempos de procesamiento tienden a aumentar o a presentar mayor dispersión a medida que crece la red, el costo computacional de ANTop se mantiene prácticamente constante en todos los escenarios evaluados. Esta propiedad resulta particularmente relevante en sistemas satelitales, donde los recursos de procesamiento disponibles a bordo suelen ser limitados.

En términos de fiabilidad, ANTop mantiene una tasa de entrega del $100\%$ en todos los escenarios considerados, incluso bajo niveles elevados de fallas en los enlaces. Por su parte, CGR \textit{per-neighbor} logra resultados comparables en la mayoría de los casos, aunque con un costo computacional significativamente mayor. En contraste, la variante CGR \textit{one-best} presenta un comportamiento menos robusto, con reducciones apreciables en la tasa de entrega y una mayor variabilidad tanto en latencia como en costo de procesamiento.

En conjunto, estos resultados sugieren que el enfoque distribuido de ANTop permite alcanzar un equilibrio favorable entre eficiencia, robustez y estabilidad computacional en redes satelitales con niveles relativamente altos de conectividad. En los escenarios evaluados, este enfoque logra superar de forma consistente a las variantes de CGR consideradas, especialmente en términos de latencia y previsibilidad del comportamiento del sistema.

\section{Trabajos futuros}

Si bien los resultados obtenidos en este trabajo permiten caracterizar el comportamiento de ANTop y compararlo con distintas variantes de CGR, existen diversas líneas de investigación que podrían profundizar y extender los hallazgos presentados.

En primer lugar, los escenarios evaluados en este trabajo corresponden a constelaciones satelitales con niveles relativamente altos de conectividad, donde múltiples rutas potenciales suelen existir entre origen y destino. Una línea natural de investigación futura consiste en analizar el desempeño de ANTop en entornos de mayor discontinuidad, donde los nodos deban esperar períodos prolongados para la aparición de nuevos enlaces. En este contexto, podría explorarse el desarrollo de mecanismos adicionales que permitan aumentar la resiliencia del algoritmo frente a largas interrupciones de conectividad, acercando su comportamiento al de enfoques basados en planes de contactos globales.

Otra posible extensión consiste en analizar alternativas al sistema de discretización espacial utilizado en este trabajo. La implementación actual se basa en la biblioteca H3, que divide la superficie terrestre en celdas hexagonales jerárquicas. Si bien esta representación ofrece ventajas importantes en términos de uniformidad de área y regularidad en la distancia entre vecinos, presenta algunas limitaciones prácticas. En particular, la resolución más baja del sistema contiene 122 celdas, lo que puede resultar excesivo para constelaciones satelitales relativamente pequeñas. Por ejemplo, el sistema GPS opera actualmente con aproximadamente 31 satélites, por lo que una discretización más flexible podría permitir representar de forma más natural ciertos escenarios.

Adicionalmente, el sistema H3 no proporciona un sistema de coordenadas global único, lo que introduce complejidades adicionales al momento de representar la topología completa de la red. En la implementación desarrollada en este trabajo fue necesario utilizar múltiples sistemas locales y estructuras auxiliares para construir una representación global coherente. En este sentido, resultaría interesante explorar otras bibliotecas o esquemas de particionamiento espacial que permitan simplificar esta representación o adaptarla mejor a las necesidades específicas de redes satelitales.

Asimismo, una línea de trabajo relevante consiste en profundizar el análisis del comportamiento de las distintas variantes de CGR en los escenarios evaluados. En particular, sería valioso investigar con mayor detalle las causas de la degradación observada en métricas como \textit{Arrival Time} y \textit{Delivery Ratio}, con el objetivo de identificar posibles limitaciones estructurales de estos algoritmos y proponer mejoras que permitan adaptar su funcionamiento a constelaciones satelitales modernas de gran escala.
